\begin{proof}[Proof of \cref{obj:thm:tail-essential-full-record-converse}]
Fix an admissible tuple \(v=(p,\alpha,\beta,\gamma)\), and put
\(w=(\alpha,\beta,\gamma)\). Throughout, use the quantities and selected
constructions of
\cref{obj:def:sharp-frontier-scales,obj:def:converse-handle}.

\begin{enumerate}
\item
Fix \(n\ge2\), and first suppose \(F(p)<1\).
By \cref{obj:lem:rough-full-record-sharp-lower}, both selected priors are
normalized, and
\[
\begin{gathered}
\operatorname{supp}_+(\pi_{0,n,v})\subseteq H_0(v),\\
\operatorname{supp}_+(\pi_{1,n,v})
\subseteq
\{P\in\mathcal M_v:c_vn^{-E_4(p)}\le d(P)<d_0\},\\
d_0\le D_v,\qquad
\operatorname{TV}(\mathbb P_{0,n,v},\mathbb P_{1,n,v})<\frac12,\qquad
B_n(v,c_vn^{-E_4(p)})\ge\frac25.
\end{gathered}
\]
These conclusions cover every \(n\ge2\), including the paired-tent
selection below \(N_{\mathrm{low}}(v)\).

The scale identification follows by subtracting the two exponents:
\[
\begin{aligned}
E_4(p)-E_0(p)
&=
\frac{2S(2\gamma+q)-2\gamma qD_p}
     {(2\gamma+q)D_p}\\
&=
\frac{2\gamma q}{(2\gamma+q)D_p}\{F(p)-1\}.
\end{aligned}
\]
Admissibility gives \(q>0\), \(S>0\), \(\gamma>0\), and \(D_p>0\).
Thus \(F(p)<1\) implies
\[
E(p)=E_4(p),\qquad \rho_n(v)=n^{-E_4(p)}.
\]

We also verify the arm support at the selected magnitude. When
\(n\ge N_{\mathrm{low}}(v)\), the rank conditions are supplied by
\cref{obj:lem:rough-full-record-sharp-lower}, and the tuning agrees with
\cref{obj:lem:copula-frame-legality}:
\[
b=\frac{K^{-\beta}}{256},\qquad
\varepsilon_{n,v}=(16b)^{1/q}
=16^{-1/q}K^{-\beta/q},\qquad
L_{n,v}=\varepsilon_{n,v}^{-1/p}.
\]
The normalized table in \cref{obj:def:marked-table} assigns its outcome
mass to \(0,-L_{n,v},L_{n,v}\). Conditioning on either treatment arm
therefore retains all its mass on those three points. When
\(n<N_{\mathrm{low}}(v)\), the selected paired-tent construction has the
same arm-support conclusion by
\cref{obj:lem:paired-tent-full-record-lower}. Consequently, in either
selection,
\[
Q_{a,P}\bigl(\{0,-L_{n,v},L_{n,v}\}\mid x\bigr)=1
\]
for each \(a\in\{0,1\}\), for \(\lambda\)-almost every \(x\), and for every
positive-weight law in either prior.

Finally, normalization and nonnegative weights ensure that the
alternative prior has a positive-weight law. Choosing such a law \(P\)
gives
\[
c_v\rho_n(v)\le d(P)<d_0\le D_v.
\]
Moreover, \(N_{\mathrm{low}}(v)\ge2\), \(k_v>0\), and
\(c_0^{\mathrm e}>0\), so the defining minimum for \(c_v\) is positive.
Since \(n>0\), also \(\rho_n(v)>0\). Hence
\[
0<c_v\rho_n(v)<D_v,
\]
and every positive-weight alternative law has distance strictly below
\(D_v\). This proves the first group of conclusions in the rough phase.
% lean: CausalSmith.Stat.FinitepHomogeneityDensegamma.rough_phase_converse_assembly

\item
Suppose instead \(F(p)\ge1\). The exponent-difference identity from the
preceding step gives
\[
E(p)=E_0(p),\qquad
c_v=c_0^{\mathrm e},\qquad
c_v\rho_n(v)=c_0^{\mathrm e}n^{-E_0(p)}.
\]
The selector of \cref{obj:def:converse-handle} chooses the paired-tent
priors for every \(n\ge2\). Their finite sign distributions are
normalized. All null sign configurations give the same zero-effect
null law \(P_0\), so
\[
\pi_{0,n,v}=\delta_{P_0},\qquad
\mathbb P_{0,n,v}=P_0^{\otimes n}.
\]
By \cref{obj:lem:paired-tent-full-record-lower},
\[
\begin{gathered}
P_0\in H_0(v),\\
\operatorname{supp}_+(\pi_{1,n,v})
\subseteq
\{P\in\mathcal M_v:
c_0^{\mathrm e}n^{-E_0(p)}\le d(P)<d_0\},\\
d_0\le D_v,\qquad
\operatorname{TV}(\mathbb P_{0,n,v},\mathbb P_{1,n,v})<\frac12,\\
B_n(v,c_0^{\mathrm e}n^{-E_0(p)})\ge\frac25.
\end{gathered}
\]
That construction also gives the required three-point conditional arm
support at \(L_{n,v}\). Its alternative prior has a positive-weight
law; for any such law,
\[
c_v\rho_n(v)\le d(P)<d_0\le D_v.
\]
Positivity of \(c_0^{\mathrm e}\) and \(n^{-E_0(p)}\) yields
\(0<c_v\rho_n(v)<D_v\). Substituting the displayed scale identity proves
the first group of conclusions in this phase as well. The equality
case \(F(p)=1\) is included here.
% lean: CausalSmith.Stat.FinitepHomogeneityDensegamma.nonrough_phase_converse_assembly

\item
We next prove divergence of the selected magnitude, again splitting at
\(F(p)<1\). In that phase, for all
\(n\ge N_{\mathrm{low}}(v)\), the selected fine rank satisfies
\[
K\ge H_{\mathrm{fine}}n^{2/D_p}\longrightarrow\infty,
\]
and moment normalization gives
\[
L_{n,v}
=
\left(16^{-1/q}K^{-\beta/q}\right)^{-1/p}
=
16^{1/(pq)}K^{\beta/(pq)}.
\]
Both \(2/D_p\) and \(\beta/(pq)\) are positive, so
\(L_{n,v}\to\infty\). The finitely many preceding selections leave this
limit unchanged.

When \(F(p)\ge1\), the paired-tent formula applies throughout:
\[
N=2\left\lceil n^{2q/(2\gamma+q)}\right\rceil
\ge n^{2q/(2\gamma+q)}\longrightarrow\infty,
\qquad
L_{n,v}=N^{\gamma/(p-1)}.
\]
Here \(2q/(2\gamma+q)>0\) and \(\gamma/(p-1)>0\), proving the same
divergence.
% lean: CausalSmith.Stat.FinitepHomogeneityDensegamma.magnitude_tendsto_atTop

\item
Assume \(p<2\). At the matched smoothness tuple, write
\[
D_2=1+2S+\frac{S}{2\gamma}.
\]
The identity \(q^{-1}=2+(2-p)/(p-1)\) gives
\[
D_p=D_2+\frac{\beta(2-p)}{p-1}.
\]
Direct subtraction then yields the two strictly positive differences
\[
E_0(2)-E_0(p)
=
\frac{4\gamma^2(2-p)}
     {(4\gamma+1)(2\gamma p+p-1)}>0,
\]
\[
E_4(2)-E_4(p)
=
\frac{2S\beta(2-p)}
     {(p-1)D_pD_2}>0.
\]
All denominators are positive by admissibility. Since \(E(p)\) is the
minimum of its two component exponents,
\[
E(p)\le E_0(p)<E_0(2),\qquad
E(p)\le E_4(p)<E_4(2).
\]
Thus the exponent gap, defined on this face by
\[
\delta_v=E(2)-E(p),
\]
is strictly positive. For every positive integer \(n\),
\[
\frac{\rho_n(v)}{\rho_n^{\mathrm b}(w)}
=n^{\delta_v}\longrightarrow\infty.
\]
% lean: CausalSmith.Stat.FinitepHomogeneityDensegamma.matched_scale_ratio_tendsto

\item
Continue with \(p<2\). Since \(c_v>0\), the preceding limit permits an
integer \(N_{\mathrm{dom}}(v)\in\mathbb N\) such that
\[
\frac{\rho_n(v)}{\rho_n^{\mathrm b}(w)}
\ge\frac{C_w^{\mathrm b}}{c_v}
\qquad
\text{for every integer }n\ge N_{\mathrm{dom}}(v)\text{ with }n>0.
\]
Multiplying by the positive bounded scale and by \(c_v\) gives
\[
C_w^{\mathrm b}\rho_n^{\mathrm b}(w)\le c_v\rho_n(v).
\]
Define the required threshold by
\[
N(v)=\max\{2,N_{\mathrm{dom}}(v),N_w^{\mathrm b}\}.
\]
For \(n\ge N(v)\), \cref{obj:thm:original-record-attainable-rate}, applied
at \((2,w)\), and \cref{obj:prop:bounded-submodel-comparison} give
\[
C_w^{\mathrm b}\rho_n^{\mathrm b}(w)<d_0\le D_w^{\mathrm b}.
\]
The bounded total test therefore satisfies
\[
\mathsf R_n(P,\phi^{*,\mathrm b}_{n,w})\le0.0075
\quad(P\in H_0^{\mathrm b}(w)),
\]
and
\[
1-\mathsf R_n(P,\phi^{*,\mathrm b}_{n,w})\le0.0075
\quad
(P\in\mathcal M_w^{\mathrm b},\ d(P)\ge c_v\rho_n(v)).
\]
Its size is below \(1/10\), so
\[
B_n^{\mathrm b}(w,c_v\rho_n(v))\le0.0075.
\]

Let \(\pi_0,\pi_1\) be any normalized finite priors satisfying the bounded
support conditions in the statement. We establish the two-prior
inequality used to compare their mixtures. For any randomized test
\(\phi\), define its seed average on the complete record space by
\[
f_\phi:
\bigl([0,1]\times\{0,1\}\times\mathbb R\bigr)^n\longrightarrow[0,1],
\qquad
f_\phi(\mathcal D_n)=\int_{[0,1]}\phi(\mathcal D_n,U)\,\lambda(dU).
\]
Bounded product integration and finite-sum linearity give, for either
prior,
\[
\int f_\phi\,d\mathbb P_{\pi_\nu,n}
=
\int \mathsf R_n(P,\phi)\,\pi_\nu(dP),
\qquad \nu\in\{0,1\}.
\]
The bounded-function form of the total-variation inequality gives
\[
\left|
\int f_\phi\,d\mathbb P_{\pi_1,n}
-
\int f_\phi\,d\mathbb P_{\pi_0,n}
\right|
\le
\operatorname{TV}(\mathbb P_{\pi_0,n},\mathbb P_{\pi_1,n}).
\]
If \(\phi\) has size at most \(1/10\) on the bounded null, normalization
and the null support condition imply
\[
\int \mathsf R_n(P,\phi)\,\pi_0(dP)\le\frac1{10}.
\]
The alternative support condition and normalization consequently yield
\[
\begin{aligned}
\sup_{\substack{P\in\mathcal M_w^{\mathrm b}\\
                 d(P)\ge c_v\rho_n(v)}}
\{1-\mathsf R_n(P,\phi)\}
&\ge
\int \{1-\mathsf R_n(P,\phi)\}\,\pi_1(dP)\\
&\ge
\frac9{10}
-\operatorname{TV}(\mathbb P_{\pi_0,n},\mathbb P_{\pi_1,n}).
\end{aligned}
\]
Weights equal to zero contribute zero in these averages; every
positive-weight law satisfies the stipulated support condition.
Normalization ensures a positive-weight alternative law, and the zero
rejection map is an admissible test. Taking the infimum over
null-sized tests therefore gives
\[
\frac9{10}
-\operatorname{TV}(\mathbb P_{\pi_0,n},\mathbb P_{\pi_1,n})
\le B_n^{\mathrm b}(w,c_v\rho_n(v)).
\]
Combining this with the attained risk bound proves
\[
\operatorname{TV}(\mathbb P_{\pi_0,n},\mathbb P_{\pi_1,n})
\ge0.8925\ge\frac12.
\]
The threshold depends on \(v\), and the argument applies to every
normalized finite prior pair in the statement.
% lean: CausalSmith.Stat.FinitepHomogeneityDensegamma.bounded_tail_essentiality

\item
Now assume \(p=2\), and fix \(n\ge2\). Then
\(v=(2,\alpha,\beta,\gamma)\).
We first verify the original-model support conclusions for the selected
bounded priors, splitting at \(F(2)<1\).

In the rough phase, \cref{obj:lem:rough-full-record-sharp-lower} supplies
\[
\begin{gathered}
\operatorname{supp}_+(\pi^{\mathrm b}_{0,n,v})
\subseteq H_0^{\mathrm b}(w),\\
\operatorname{supp}_+(\pi^{\mathrm b}_{1,n,v})
\subseteq
\{P\in\mathcal M_w^{\mathrm b}:
c_vn^{-E_4(2)}\le d(P)<d_0\}.
\end{gathered}
\]
By the phase identity, \(\rho_n(v)=n^{-E_4(2)}\).
The model inclusion of \cref{obj:prop:bounded-submodel-comparison}
preserves the effect, and the null constancy condition is the same in
both models. Hence
\[
\begin{gathered}
\operatorname{supp}_+(\pi^{\mathrm b}_{0,n,v})\subseteq H_0(v),\\
\operatorname{supp}_+(\pi^{\mathrm b}_{1,n,v})
\subseteq
\{P\in\mathcal M_v:c_v\rho_n(v)\le d(P)\}.
\end{gathered}
\]

In the complementary phase \(F(2)\ge1\), the selected bounded priors are
the signed-binary paired-tent priors. To identify their null support,
define the balanced deterministic zero-effect law by
\[
P_{\mathrm{zero}}(dx,dA,dy)
=
\lambda(dx)\,
\frac{\delta_0+\delta_1}{2}(dA)\,
\delta_0(dy).
\]
It has propensity \(1/2\), baseline mean zero, and effect zero. Its
smoothness conditions hold because these functions are constant, and
its raw moment is zero. Its signed-binary conversion from
\cref{obj:prop:bounded-submodel-comparison} is
\[
P_{\mathrm{zero}}^{\mathrm{bin}}(dx,dA,dy)
=
\lambda(dx)\,
\frac{\delta_0+\delta_1}{2}(dA)\,
\frac{\delta_{-1}+\delta_1}{2}(dy).
\]
This conversion preserves the means and effect, and its conditional
second moment is one. Thus
\(P_{\mathrm{zero}}^{\mathrm{bin}}\in H_0(v)\).
Every selected null sign configuration gives precisely this law, so
\[
\operatorname{supp}_+(\pi^{\mathrm b}_{0,n,v})
\subseteq\{P_{\mathrm{zero}}^{\mathrm{bin}}\}\subseteq H_0(v).
\]
For the alternatives,
\cref{obj:lem:paired-tent-full-record-lower} and
\cref{obj:prop:bounded-submodel-comparison} give original-model
membership and
\[
d(P)\ge c_0^{\mathrm e}n^{-E_0(2)}
=c_v\rho_n(v)
\qquad
(P\in\operatorname{supp}_+(\pi^{\mathrm b}_{1,n,v})).
\]
This proves the original-model support conclusions in both phases.
% lean: CausalSmith.Stat.FinitepHomogeneityDensegamma.second_moment_rough_prior_support

\item
The bounded risk lower bound follows from the same phase split.
When \(F(2)<1\),
\cref{obj:lem:rough-full-record-sharp-lower} gives
\[
B_n^{\mathrm b}(w,c_vn^{-E_4(2)})\ge\frac25.
\]
When \(F(2)\ge1\),
\cref{obj:lem:paired-tent-full-record-lower} gives
\[
B_n^{\mathrm b}(w,c_0^{\mathrm e}n^{-E_0(2)})\ge\frac25.
\]
The phase identities identify each displayed separation with
\(c_v\rho_n(v)\). The risk comparison in
\cref{obj:prop:bounded-submodel-comparison} consequently yields
\[
\frac25
\le B_n^{\mathrm b}(w,c_v\rho_n(v))
\le B_n(v,c_v\rho_n(v)).
\]
% lean: CausalSmith.Stat.FinitepHomogeneityDensegamma.second_moment_testingRisk_lower

\item
We assemble the remaining second-moment conclusions. Evaluation at
\(v=(2,w)\), together with
\cref{obj:prop:bounded-submodel-comparison}, gives
\[
\begin{gathered}
\rho_n(v)=\rho_n^{\mathrm b}(w),\qquad
c_v=c_w^{\mathrm b},\qquad
C_v=C_w^{\mathrm b},\qquad
N_v=N_w^{\mathrm b},\\
D_v=D_w^{\mathrm b}.
\end{gathered}
\]
For \(F(2)<1\), the selected-prior conclusions of
\cref{obj:lem:rough-full-record-sharp-lower} apply for every \(n\ge2\),
covering the bounded rough construction above its threshold and the
signed-binary paired-tent construction below it. For \(F(2)\ge1\), use
the signed-binary conclusions of
\cref{obj:lem:paired-tent-full-record-lower}. The signed-binary model
inclusion of \cref{obj:prop:bounded-submodel-comparison} supplies bounded
membership, with null constancy retained. In either case, both selected
bounded priors are normalized and
\[
\begin{gathered}
\operatorname{supp}_+(\pi^{\mathrm b}_{0,n,v})
\subseteq H_0^{\mathrm b}(w),\\
\operatorname{supp}_+(\pi^{\mathrm b}_{1,n,v})
\subseteq
\{P\in\mathcal M_w^{\mathrm b}:
c_v\rho_n(v)\le d(P)<d_0\},\\
\operatorname{TV}(\mathbb P^{\mathrm b}_{0,n,v},
                  \mathbb P^{\mathrm b}_{1,n,v})<\frac12,\qquad
B_n^{\mathrm b}(w,c_v\rho_n(v))\ge\frac25.
\end{gathered}
\]
For any finite representation of either prior, nonnegative weights
summing to one satisfy
\[
\sum_{j:\,\omega_j>0}\omega_j=\sum_j\omega_j=1.
\]
Indeed, every omitted weight is zero. At least one weight is positive,
and combining weights of coincident laws preserves their sum.
Therefore
\[
\operatorname{supp}_+(\pi^{\mathrm b}_{\nu,n,v})\ne\varnothing,
\qquad
\sum_{P\in\operatorname{supp}_+(\pi^{\mathrm b}_{\nu,n,v})}
\pi^{\mathrm b}_{\nu,n,v}(\{P\})=1,
\qquad \nu\in\{0,1\}.
\]
Choose a positive-weight alternative law \(P\). Its placement, and the
witness lower bound from
\cref{obj:prop:bounded-submodel-comparison}, give
\[
0<c_v\rho_n(v)\le d(P)<d_0\le D_w^{\mathrm b}.
\]
% lean: CausalSmith.Stat.FinitepHomogeneityDensegamma.second_moment_bounded_receipt

\item
We obtain the lower capped-radius bound first in the bounded model.
Every successful bounded separation satisfies
\[
r_{\mathrm{succ}}\in(0,D_w^{\mathrm b}),
\qquad
B_n^{\mathrm b}(w,r_{\mathrm{succ}})\le\frac1{10}.
\]
If \(r_{\mathrm{succ}}<c_v\rho_n(v)\), containment of the separated
alternative classes gives
\[
B_n^{\mathrm b}(w,c_v\rho_n(v))
\le B_n^{\mathrm b}(w,r_{\mathrm{succ}})
\le\frac1{10},
\]
contradicting the lower bound \(2/5\). The cap
\(D_w^{\mathrm b}\) also exceeds \(c_v\rho_n(v)\). Thus every member of
the nonempty set defining the capped infimum is at least
\(c_v\rho_n(v)\), and
\[
c_v\rho_n(v)\le r_n^{*,\mathrm b}(w).
\]
Now \cref{obj:prop:bounded-submodel-comparison} gives
\[
c_v\rho_n(v)
\le r_n^{*,\mathrm b}(w)
=r_n^{*,\mathrm{bin}}(w)
\le r_n^*(v).
\]
The upper capped-radius bounds and both size bounds follow from
\cref{obj:thm:original-record-attainable-rate}:
\[
r_n^*(v)\le C_v\rho_n(v),\qquad
r_n^{*,\mathrm b}(w)\le C_w^{\mathrm b}\rho_n^{\mathrm b}(w),
\]
\[
\mathsf R_n(P,\phi^*_{n,v})\le0.0075
\quad(P\in H_0(v)),\qquad
\mathsf R_n(P,\phi^{*,\mathrm b}_{n,w})\le0.0075
\quad(P\in H_0^{\mathrm b}(w)).
\]
Substituting the matching scale and multiplier identities proves both
asserted two-sided radius bounds.
% lean: CausalSmith.Stat.FinitepHomogeneityDensegamma.second_moment_capped_lower

\item
Suppose \(C_v\rho_n(v)<D_v\). The matching identities give
\[
C_w^{\mathrm b}\rho_n^{\mathrm b}(w)<D_w^{\mathrm b}.
\]
To prove nonemptiness in the bounded model, suppose its separated class
were empty. Every bounded-model law would then satisfy
\[
d(P)<C_w^{\mathrm b}\rho_n^{\mathrm b}(w).
\]
The bounded model is nonempty, as witnessed by a positive-weight law
from the selected alternative prior. The displayed bound would imply
\[
D_w^{\mathrm b}
\le C_w^{\mathrm b}\rho_n^{\mathrm b}(w),
\]
contradicting the strict inequality. Hence there is a bounded-model law
with
\[
d(P)\ge C_w^{\mathrm b}\rho_n^{\mathrm b}(w)=C_v\rho_n(v).
\]
Its original-model membership follows from
\cref{obj:prop:bounded-submodel-comparison}, proving nonemptiness of both
separated classes.

Under this same strict-separation condition,
\cref{obj:thm:original-record-attainable-rate} supplies
\[
1-\mathsf R_n(P,\phi^*_{n,v})\le0.0075
\quad(P\in\mathcal M_v,\ d(P)\ge C_v\rho_n(v)),
\]
and
\[
1-\mathsf R_n(P,\phi^{*,\mathrm b}_{n,w})\le0.0075
\quad
(P\in\mathcal M_w^{\mathrm b},\
d(P)\ge C_w^{\mathrm b}\rho_n^{\mathrm b}(w)).
\]
% lean: CausalSmith.Stat.FinitepHomogeneityDensegamma.second_moment_face_package

\item
Finally, let \(n\ge N_v\). Positivity of the component exponents implies
\(E(2)>0\), and
\cref{obj:thm:original-record-attainable-rate} gives \(C_v>0\).
The threshold formula yields
\[
n\ge
\max\left\{4,\left(\frac{2C_v}{d_0}\right)^{1/E(2)}\right\}.
\]
Since raising positive numbers to the power \(-E(2)\) reverses their
order,
\[
\begin{aligned}
C_v\rho_n(v)
&=C_vn^{-E(2)}\\
&\le
C_v\left\{
\left(\frac{2C_v}{d_0}\right)^{1/E(2)}
\right\}^{-E(2)}
=\frac{d_0}{2}.
\end{aligned}
\]
Also, the witness-threshold conclusion of
\cref{obj:thm:original-record-attainable-rate} gives
\[
C_v\rho_n(v)<d_0\le D_v.
\]
The matching identities therefore imply
\[
C_w^{\mathrm b}\rho_n^{\mathrm b}(w)<D_w^{\mathrm b}.
\]
At \(n=2,3\), the small-sample branch of
\cref{obj:def:explicit-score-ledger} sets the total rule to zero.
Evaluating that same branch at \(v=(2,w)\) gives, for every sample \(z\),
\[
\phi^*_{n,v}(z)=0,\qquad
\phi^{*,\mathrm b}_{n,w}(z)=0.
\]
This completes all conclusions on the second-moment face.
% lean: CausalSmith.Stat.FinitepHomogeneityDensegamma.second_moment_attainment_half
\end{enumerate}
\end{proof}
